\section{Limitations}
\label{sec:limitations}

The results should be interpreted within four boundaries. First, the held-out evaluation contains 30,511 stock-quarter observations but only eight formation quarters from 2024~Q1 through 2025~Q4 and 502 trading days in the continuous portfolio analysis. The cross-section is therefore large relative to the time dimension. Chronological separation protects the reported test from model selection, but a single short test period cannot establish stability across market regimes.

Second, Form 13F is a partial and delayed representation of institutional portfolios. It reports specified positions quarterly for managers meeting the regulatory threshold \parencite{sec2023form13f}; it does not reveal the complete investor base, short positions, or trading between reporting dates. The ownership and network variables consequently measure reported institutional long-equity exposure available at the common information cutoff, not the full set of positions or the trading process that generated them. The analysis is predictive, and neither the fitted relationships nor the SHAP diagnostics identify causal effects of ownership structure on later drawdowns.

Third, the model comparisons are conditional on the implemented specifications and search ranges. Ridge uses the prespecified penalty $\alpha=1$, whereas XGBoost and the selected graph architecture undergo controlled validation refinements. Ridge and XGBoost minimize squared-error objectives, while the graph models use $L_1$ loss. Hypotheses~H3 and~H4 should therefore be interpreted as comparisons among the implemented specifications rather than exhaustive comparisons among optimally tuned model classes. Manager similarity is represented through a 25-nearest-neighbour graph, Louvain communities are estimated separately each quarter, and the tabular models receive the network features retained by the prespecified redundancy screen. The graph comparison covers weighted GraphSAGE, GAT, and Temporal GraphSAGE under a bounded tuning exercise. The absence of consistent incremental performance applies to these representations; it does not rule out alternative edge definitions, community methods, temporal structures, or graph architectures. The comparisons are also based on point estimates rather than formal tests of differences in predictive accuracy. Given the small gaps among the finalists and the limited time dimension described above, the evidence does not establish that these differences are statistically distinguishable; a longer test period and inference that accounts for cross-sectional and temporal dependence would strengthen the comparison. In addition, each graph architecture is estimated using a single fixed random seed, so sensitivity to random initialization is not measured. These qualifications do not change the validation-stage selection decision: XGBoost is retained because it combines similar MAE and ranking performance with stronger RMSE and $R^2$ and substantially lower complexity, rather than because it achieves the lowest MAE. The similar metric pattern observed in validation and test provides some cross-period reassurance, but it does not resolve the inferential or seed-related uncertainty.

Fourth, the portfolio tests do not constitute a complete implementation study. The strongest long-only result is obtained under equal weighting and becomes negligible under value weighting, indicating sensitivity to the weighting rule and the influence of smaller stocks. The CRSP total-market index is an external reference rather than a matched benchmark. Transaction costs are modeled through fixed turnover-based assumptions, while market impact, capacity, taxes, stock-borrow availability, borrowing fees, and financing costs are not fully incorporated. The long--short and volatility-managed results should therefore be interpreted as evidence of economic separation, not as verified net-of-cost investment performance.
